% Chapter 7: Conclusion and Future Work

\chapter{Conclusion and Future Work}
\label{chap:conclusion}

\section{Summary of Contributions}
\label{sec:conclusion-contributions}

This dissertation contributed a comparable experimental platform for UAV-to-GCS authentication, a permissioned-blockchain identity registry with explicit smart-contract lifecycle rules, an X.509 PKI baseline under identical challenge--response semantics, and a reproducible 53-step evaluation matrix with exported evidence (security/audit: \texttt{dissertation\_final\_20260831T015600Z}; Method~B scale/ABBA: \texttt{method\_b\_20260904}).

The hardened evaluation demonstrated that both mechanisms can achieve full normal-authentication acceptance (750/750 measured repetitions each) and zero false acceptances in the adversarial battery, while exhibiting materially different performance characteristics. On the Method~B scale re-run, X.509 provided lower latency (approximately 2.03~ms mean at $n=500$, $c=1$) and roughly twice the completed throughput of blockchain under high concurrency (76.98 versus 38.66~rps at $n=500$, $c=50$).

\section{Answers to the Research Questions}
\label{sec:conclusion-rqs}

\textbf{RQ1.} Both implemented mechanisms resisted the executed spoofing, replay, tampering, unauthorised-access and revocation scenarios without false acceptances, subject to the stated threat boundary and local deployment assumptions.

\textbf{RQ2.} Blockchain imposed higher decision latency and CPU overhead than X.509; optional audit transactions added sub-second submission and $\approx$ 2~s confirmation latency with measurable gas cost.

\textbf{RQ3.} X.509 reached 76.98~rps at $n=500$, $c=50$ on the evaluation host with stable low-concurrency latency up to 500 identities. Blockchain throughput at the same point was 38.66~rps with pronounced latency growth under concurrent registry lookups.

\textbf{RQ4.} Blockchain is conditionally justified when append-only auditability and shared registry semantics outweigh latency and infrastructure complexity; X.509 remains preferable for latency-critical, centrally administered deployments.

\section{Future Work}
\label{sec:conclusion-future}

Recommended extensions include: (i)~wide-area and lossy-channel emulation for authentication latency; (ii)~validator fault-injection and quorum-loss tests on the QBFT network; (iii)~persistent nonce storage and GCS failover behaviour; (iv)~formal verification or third-party audit of the registry contract; (v)~hardware-backed key storage on representative UAV platforms; and (vi)~integration with transport-layer security and mission-system authorisation policies beyond the authentication boundary defined in Chapter~\ref{chap:research-design-methodology}.

\section{Closing Remarks}
\label{sec:conclusion-closing}

Permissioned blockchain technology can support UAV identity management and auditable authentication decisions, but it is not a drop-in replacement for PKIX on performance grounds alone. The experimental evidence from this study supports mechanism selection based on explicit operational priorities---latency, auditability, administrative model and infrastructure cost---rather than on architectural preference unsupported by measurement.
